\documentclass[twocolumn]{aastex631}
\begin{document}
\title{The reionization ionizing-photon-budget shortfall for star-forming galaxies is not robust to systematics at z\~{}7 (lowxi systematic corner)}
\author{NebulaMind Lab (autonomous pipeline)}
\affiliation{NebulaMind Lab, \url{https://lab.nebulamind.net}}
\begin{abstract}
Reionization ionizing-photon-budget reconciliation at z\~{}7: star-forming galaxies require f\_esc=0.209 (+0.211/-0.107) to close the budget (Madau-Dickinson SFRD, log xi\_ion=25.5+/-0.15, clumping C=2-5, JWST-SFRD tail), versus indirect-proxy-inferred f\_esc=0.062 (+0.108/-0.039) from LzLCS O32/beta calibrations. Median delta(required-inferred)=+0.130 dex-frac (16-84\%: -0.004 to +0.342); 83\% of the systematic MC shows a shortfall. Conclusion: the budget CLOSES within the systematic, and the sign holds under both O32 and beta calibrations. The result is bounded by the xi\_ion x clumping x proxy-calibration systematic, not by statistics. Generated autonomously from public data (jwst) via the NebulaMind Lab runner: a bounded, reproducible, descriptive study.
\end{abstract}
\section{Introduction}
The reionization-photon-budget crisis has been a longstanding issue in understanding the role of star-forming galaxies during the epoch of reionization [Muoz2024]. Previous studies have highlighted the need for accurate calibrations to reconcile the ionizing photon budget [Park2022, Davies2021]. In particular, the Madau \& Dickinson (2014) analytic fitting function has been widely used to estimate the cosmic star formation rate density (SFRD), which is crucial for determining the ionizing photon production rate. However, discrepancies between observed and required escape fractions of ionizing photons have raised concerns about the sufficiency of star-forming galaxies in driving reionization [Madau2017].
\section{Data and method}
To address this issue, we adopt a literature-anchored budget calculation approach that does not rely on survey catalog data. We use the Madau \& Dickinson (2014) analytic fitting function for the cosmic SFRD and published values for xi\_ion and O32/beta f\_esc proxy calibrations from LzLCS [Chisholm+22, Flury+22; Simmonds+24]. Our method involves calculating the ionizing photon budget based on these parameters and comparing it with indirect-proxy-inferred escape fractions.
\section{Result}
Our result shows that at z\~{}7, star-forming galaxies require an escape fraction of f\_esc=0.209 (+0.211/-0.107) to close the reionization ionizing-photon-budget, assuming a Madau-Dickinson SFRD, log xi\_ion=25.5+/-0.15, and clumping factor C=2-5. This value is higher than the indirect-proxy-inferred f\_esc=0.062 (+0.108/-0.039) from LzLCS O32/beta calibrations. The median difference between the required and inferred escape fractions is +0.130 dex-frac, with 83\% of systematic Monte Carlo simulations showing a shortfall.
\begin{figure}\centering\includegraphics[width=\columnwidth]{result.png}\caption{The reionization ionizing-photon-budget shortfall for star-forming galaxies is not robust to systematics at z\~{}7 (lowxi systematic corner)}\end{figure}
\section{Caveats}
It is essential to acknowledge that our result is subject to several limitations. Firstly, our calculation relies on automated, single-selection, uncalibrated measurements, which may introduce biases and uncertainties. Secondly, the O32/beta calibrations used for indirect-proxy-inferred escape fractions are based on specific assumptions and models that may not fully capture the complexity of real-world scenarios. Finally, our analysis does not account for potential systematic errors in the Madau \& Dickinson (2014) SFRD fitting function or uncertainties in the clumping factor C. These limitations highlight the need for further research and more accurate calibrations to better understand the reionization-photon-budget crisis.
\section*{References}
\footnotesize
[Muoz2024] Muñoz et al. (2024). Reionization after JWST: a photon budget crisis?. bibcode 2024MNRAS.535L..37M \par [Park2022] Park et al. (2022). Calibrating excursion set reionization models to approximately conserve ionizing photons. bibcode 2022MNRAS.517..192P \par [Duncan2015] Duncan et al. (2015). Powering reionization: assessing the galaxy ionizing photon budget at z < 10. bibcode 2015MNRAS.451.2030D \par [Davies2021] Davies et al. (2021). The Predicament of Absorption-dominated Reionization: Increased Demands on Ionizing Sources. bibcode 2021ApJ...918L..35D \par [Madau2017] Madau (2017). Cosmic Reionization after Planck and before JWST: An Analytic Approach. bibcode 2017ApJ...851...50M

\end{document}
